(1) 一次腔内往返的几何长度为 $2L$、光程为 $2nL$，对应相位 $\delta=(2\pi/\lambda)(2nL)=4\pi nL/\lambda$。频率表示为 $\delta=4\pi nL\nu/c$，暂假设 $n$ 在所考察线宽内无明显色散，且两镜的反射相位已统一纳入共振相位约定。
(2) 每次额外往返产生复振幅因子 $Re^{i\delta}$，因为两次振幅反射率的乘积模为 $R$。最初透射振幅系数模为 $1-R$，所以归一化复透射幅度为 $(1-R)/(1-Re^{i\delta})$，由此 $\mathcal T=\dfrac{(1-R)^2}{1+R^2-2R\cos\delta}=\left[1+\dfrac{4R}{(1-R)^2}\sin^2(\delta/2)\right]^{-1}$。这利用了无吸收对称镜，不适用于镜透过率未知的有损腔。
(3) 当 $\delta=2\pi m$，所有往返光束同相，透过率达到 $\mathcal T=1$；腔的纵模频率为 $\nu_m=mc/(2nL)$（常数反射相位可改变整数标记但不改间隔）。因此相邻共振峰的自由谱距是 $\Delta\nu_{\rm FSR}=c/(2nL)$。只有两镜对称且无损时在共振处透过率才恰好达到一。
(4) 记 $F=4R/(1-R)^2$，在共振峰附近 $\sin^2(\delta/2)\simeq\delta^2/4$，故半高两侧相移大小为 $|\delta_{1/2}|\simeq2/\sqrt F=(1-R)/\sqrt R$。相位半高全宽为 $\Delta\delta_{\rm FWHM}\simeq2(1-R)/\sqrt R$，与自由谱距 $2\pi$ 比得精细度 $\mathcal F\simeq\pi\sqrt R/(1-R)$。当 $R\to0$ 时峰之间已无明显分辨，窄线宽近似不成立。
